\section*{第 \ref{ch:b2:climate-proxies} 章——气候档案与代用指标}

\begin{solution}{exo:b2:climate-proxies:1}
(a) $\delta = (27.5 \times (-30) + 1350 \times 0)/(1350 + 27.5) = -825/1377.5 = -0.60$ ‰。
(b) 若海洋保有全部水量，则为 $-0.60$ ‰，比今天轻
\num{0.60} ‰。(c) $27.5/1377.5 = \qty{2.0}{\percent}$。
\end{solution}

\begin{solution}{exo:b2:climate-proxies:2}
(a) $T = 16.1 - 4.76\,(2.50 - (0 - 0.27)) = 16.1 - 4.76 \times 2.77 = 16.1 - 13.19 =
\qty{2.9}{\celsius}$。(b) $T = 16.1 - 4.76\,(2.50 - 0.73) = 16.1 - 8.43 = \qty{7.7}{\celsius}$：同一
壳体在水体更重时给出的温度高出近 5 度，这正是
必须知道水体组成的原因。(c) 壳体偏重 $2.71 - 2.50 = 0.21$ ‰，
故其水温低 $4.76 \times 0.21 = \qty{1.0}{\celsius}$，为 \qty{1.9}{\celsius}。
\end{solution}

\begin{solution}{exo:b2:climate-proxies:3}
(a) $0.4\,\mathrm{e}^{0.9} = \qty{0.98}{mmol/mol}$ 与 $0.4\,\mathrm{e}^{1.8} =
\qty{2.42}{mmol/mol}$。(b) 因子为 $\mathrm{e}^{0.9} = 2.46$；每度 $\mathrm{e}^{0.09}
= 1.094$，每度增长 \qty{9.4}{\percent}。(c) $T = \ln(\text{Mg/Ca}/0.4)/0.09$：
$\ln 5/0.09 = \qty{17.9}{\celsius}$，$\ln 2.5/0.09 = \qty{10.2}{\celsius}$。
\end{solution}

\begin{solution}{exo:b2:climate-proxies:4}
(a) $2668 - 2213 + 1 = 456$ 年。(b) 单棵树可能在局部缺失一圈，或出现
一圈无法辨读的年轮；将多棵树宽窄年轮的模式相互对照，可以
发现缺失的年轮并确定每一圈的年份，树木之间的吻合程度则衡量
定年的可靠性。(c) 在林线处，生长季的寒冷这一因子
制约所有树木，因此它们共享相同的好年与坏年模式；在温和湿润的
低地，没有单一因子占主导，树木之间的模式各不相同，\omterm{def:b2:climate-proxies:dendro}{交叉定年}可能失败。
\end{solution}

\begin{solution}{exo:b2:climate-proxies:5}
(a) $125 \times 40/(3682.2 - 125) = 5000/3557.2 = 1.41$ ‰；$134 \times 40/3548.2 =
5360/3548.2 = 1.51$ ‰。(b) 取 $H - h = 3552.2$：$3900/3552.2 = 1.10$、$4550/3552.2 = 1.28$
与 $6500/3552.2 = 1.83$ ‰。(c) $\delta_{\mathrm{ice}} = -1.0 \times 3552.2/130 = -27.3$
‰，比现今任何冰盖都重（格陵兰 $-30$ 至 $-35$，北美冰盖
$-40$，南极冰低至 $-50$）。依据真实冰量算出的质量平衡与孔隙水
测量结果不一致；约半个 ‰ 的差距是该方法真实存在的不确定度。
\end{solution}

\begin{solution}{exo:b2:climate-proxies:6}
(a) $0.372\,\mathrm{e}^{0.0936T} = 0.4\,\mathrm{e}^{0.09T}$ 给出 $\mathrm{e}^{0.0036T} =
0.4/0.372 = 1.0753$，因此 $T = \ln 1.0753/0.0036 = \qty{20.2}{\celsius}$。(b) 在
\qty{2}{\celsius} 处：$0.4\,\mathrm{e}^{0.18} = 0.479$ 与 $0.372\,\mathrm{e}^{0.1872} = 0.449$，故
拟合比值为 $0.449/0.479 = 0.937$，低 \qty{6.3}{\percent}。(c) $T = \ln(0.449/0.4)/0.09 =
\ln 1.1225/0.09 = \qty{1.3}{\celsius}$：预测的规律使壳体温度读数偏冷约
\qty{0.7}{\celsius}。冷水中，标定的选择对结果的影响
达到一度的量级。
\end{solution}

\begin{solution}{exo:b2:climate-proxies:7}
(a) 从全新世到 GI-1，冰层厚 \qty{112.19}{m}，历时 2989 年：$11\,219/2989 =
\qty{3.75}{cm}$ 每年；从 GI-16 到 GI-17，冰层厚 \qty{17.89}{m}，历时 1160 年：$1789/1160 =
\qty{1.54}{cm}$ 每年。(b) 冰层在
流动：每层被埋藏后向两侧扩展并变薄，而越深、越老的层
变薄的时间更长。(c) 当冰层薄至一厘米或更薄时，仪器
已无法分辨相邻两年，\omterm{def:b2:climate-proxies:age-model}{年龄模型}不得不
改用定位点或流动模型。
\end{solution}

\begin{solution}{exo:b2:climate-proxies:8}
(a) 区间内的比例为 $(1550 - 1492.45)/(1604.64 - 1492.45) = 57.55/112.19 = 0.513$，
故年龄为 $11\,703 + 0.513 \times 2989 = 11\,703 + 1533 = \num{13236}$ 年 b2k。(b) 1950 年比 2000 年早
50 年：比 1950 年早 \num{13186} 年。(c) $50 + 0.513 \times (93 - 50) = 72$ 年。
\end{solution}

\begin{solution}{exo:b2:climate-proxies:9}
(a) 每个不确定层贡献半年：$2 \times 93 = 186$ 与 $2 \times 1287 = 2574$
个不确定层。(b) $\frac12\sqrt{186} = 6.8$ 年，$\frac12\sqrt{2574} = 25.4$ 年。
(c) $93/6.8 = 14$ 与 $1287/25.4 = 51$。$\frac12 n$ 与 $\frac12\sqrt{n}$ 之比为
$\sqrt{n}$：它随不确定层数的平方根增长，因此两种规则之间的选择
随深度增加而越来越重要。
\end{solution}

\begin{solution}{exo:b2:climate-proxies:10}
(a) 在\omterm{prop:b2:isotope-geochemistry:mwl}{大气降水线}上，$\delta$D 的变化是 $\delta^{18}$O 的八倍，故
斜率为 $8 \times 0.8 = 6.4$ ‰（$\delta$D）每度。降幅为 $441.7 - 396.6 = 45.1$
‰，由此得降温 $45.1/6.4 = \qty{7.0}{\celsius}$。(b) 冰期海洋约
偏重 $8 \times 1.0 = 8$ ‰（$\delta$D）；由今日海洋形成的雪本应
同样偏轻 $-441.7 - 8 = -449.7$ ‰。降幅变为 $53.1$ ‰，降温
为 $53.1/6.4 = \qty{8.3}{\celsius}$。(c) $9.39 - 0.09 = \qty{9.30}{\celsius}$，故 $45.1/9.30 = 4.85$
‰ 每度。发表的序列还对冰盖海拔的变化作了校正，
其温度关系来自模拟水同位素的模式，而非
仅来自空间斜率。
\end{solution}

\begin{solution}{exo:b2:climate-proxies:11}
(a) 比值的对数依次为 $-0.528$、$-0.223$、$0.191$、$0.501$、$0.904$、$1.227$；其
平均值为 $0.345$，平均温度为 14。取 $S_{TT} = \sum (T - 14)^2 = 280$ 与 $S_{T\ell} =
\sum (T - 14)(\ell - 0.345) = 24.93$，$A = 24.93/280 = \qty{0.0890}{\per\celsius}$，$\ln B =
0.345 - 0.0890 \times 14 = -0.901$，故 $B = \qty{0.406}{mmol/mol}$。(b) 残差为
$+0.017$、$-0.034$、$+0.023$、$-0.023$、$+0.025$ 与 $-0.009$，占比值的百分之几，且无
明显趋势；拟合结果接近 $0.4\,\mathrm{e}^{0.09T}$。(c) $T = (\ln 1.30 + 0.901)/0.0890 = (0.262 +
0.901)/0.0890 = \qty{13.1}{\celsius}$。(d) $T = (\ln 0.45 + 0.901)/0.0890 = (-0.799 +
0.901)/0.0890 = \qty{1.2}{\celsius}$。该值低于标定所用样品中最冷的一个，
\qty{4}{\celsius}：此规律被外推，而在低温下，壳体的 Mg/Ca 比值已知会
在碳酸根离子的影响下偏离指数规律。
\end{solution}

\begin{solution}{exo:b2:climate-proxies:12}
(a) 按单位海洋面积计，现今海洋的水量（$H$，取 0）在各储库之间分配，
即 $h_i$（位于 $\delta_i$ 处）与其余海洋 $H - h$，其中 $h = \sum h_i$：$0 = \sum h_i\delta_i +
(H - h)\,\Delta\delta$。第一阶段后 $\Delta\delta = 60 \times 30/(3682.2 - 60) =
1800/3622.2 = 0.50$ ‰；第二阶段后 $(1800 + 70 \times 45)/(3682.2 - 130) =
4950/3552.2 = 1.39$ ‰。(b) $-4950/130 = -38.1$ ‰。(c) 该规则为 $1.39/130 = 0.0107$
‰ 每米，给出第一阶段 $0.50/0.0107 = \qty{47}{m}$，而非
\qty{60}{m}：第一批冰比平均值更重，因此其每一米所引起的海洋
变化较小。依据某一状态标定的线性规则，在
冰的平均组成发生变化时便会误读其他状态。
\end{solution}

\begin{solution}{pb:b2:climate-proxies:1}
\textbf{1.} $361.84 \times 10^6\ \unit{km^2} \times \qty{0.130}{km} = \qty{4.70e7}{km^3}$，约 47
百万立方千米。
\textbf{2.} $130 \times 40/(3682.2 - 130) = 5200/3552.2 = 1.46$ ‰。
\textbf{3.} $5000/3557.2 = 1.41$ 和 $5360/3548.2 = 1.51$ ‰。
\textbf{4.} $4550/3552.2 = 1.28$ 和 $6500/3552.2 = 1.83$ ‰。冰的组成
使结果改变 \num{0.55} ‰，海平面使其改变 \num{0.10} ‰：冰的影响更大。
\textbf{5.} $\delta_{\mathrm{ice}} = -1.0 \times 3552.2/130 = -27.3$ ‰。
\textbf{6.} $40/3682.2 = 0.0109$ ‰ 每米，四舍五入后与定年
海平面数据的 \num{0.011} 一致。
\textbf{7.} 在 $\delta^{18}\mathrm{O}_{\mathrm{sw}}$ 不变时，$\mathrm{d}T = -4.76\,\mathrm{d}\delta^{18}
\mathrm{O}_{\mathrm{b}}$，因此 $\mathrm{d}T = -1$ 给出 $\mathrm{d}\delta^{18}\mathrm{O}_{\mathrm{b}} =
1/4.76 = 0.21$ ‰。
\textbf{8.} $\delta^{18}\mathrm{O}_{\mathrm{b}} = (16.1 - T)/4.76 + \delta^{18}\mathrm{O}_{\mathrm{sw}}
- 0.27 = 14.1/4.76 - 0.27 = 2.96 - 0.27 = 2.69$ ‰。
\textbf{9.} $17.1/4.76 + 1.0 - 0.27 = 3.59 + 0.73 = 4.32$ ‰。变化量为 $4.32 - 2.69 =
1.63$ ‰：其中 \num{1.00} ‰ 来自海水，$3.0 \times 0.21 = 0.63$ ‰ 来自
\qty{3}{\celsius} 的降温。
\textbf{10.} $T = \ln(2.0/0.4)/0.09 = \ln 5/0.09 = \qty{17.9}{\celsius}$。
\textbf{11.} $0.18/0.1 = \qty{1.8}{\celsius}$ 的降温。
\textbf{12.} Mg/Ca 是方解石中两种元素的比值，由
钙化时的温度决定；它不涉及氧的同位素，因此
忽略了冰所改变的海水组成。
\textbf{13.} $4.92 - 3.23 = 1.69$ ‰。
\textbf{14.} 由 $\Delta = h \times 40/(H - h)$ 解出 $h$：$h = H\Delta/(\Delta + 40) = 3682.2
\times 1.69/41.69 = \qty{149}{m}$，超过海平面证据给出的 125 至 \qty{134}{m}：
变化的一部分必定来自降温。
\textbf{15.} $1.69 - 1.46 = 0.23$ ‰，且 $4.76 \times 0.23 = \qty{1.1}{\celsius}$。
\textbf{16.} $1.69 - 1.00 = 0.69$ ‰，且 $4.76 \times 0.69 = \qty{3.3}{\celsius}$（3.28）。
\textbf{17.} $0.23/0.25 = \qty{0.9}{\celsius}$，$0.69/0.25 = \qty{2.8}{\celsius}$。
\textbf{18.} 北大西洋深层水由 2--\qty{4}{\celsius} 冷却到 \qty{-1.1}{\celsius}，
降温 3.1 至 \qty{5.1}{\celsius}；南极底层水由 0--\qty{1}{\celsius} 冷却到
\qty{-1.9}{\celsius}，降温 1.9 至 \qty{2.9}{\celsius}；西北大西洋的 Mg/Ca 给出
\qty{1.8}{\celsius}。这些结果大致都在 2 至 \qty{5}{\celsius} 之间，围绕孔隙水校正给出的 2.8 至
\qty{3.3}{\celsius}，并高于
冰量平衡在 $-40$ ‰ 时给出的 \qty{1}{\celsius}。
\textbf{19.} $-396.6 - (-441.7) = 45.1$ ‰。
\textbf{20.} $8 \times 0.8 = 6.4$ ‰ 每度，因此 $45.1/6.4 = \qty{7.0}{\celsius}$。
\textbf{21.} $9.39 - 0.09 = \qty{9.30}{\celsius}$，且 $45.1/9.30 = 4.85$ ‰ 每度。
\textbf{22.} Dome C，\qty{9.3}{\celsius}，高于全球地表的 5 至 \qty{7}{\celsius}，高于
深海约 \qty{3}{\celsius}。
\textbf{23.} 今天深层水本已很冷，为 0 至 \qty{4}{\celsius}，而海水
不能冷却到远低于其凝固点的温度，末次冰期的底层水几乎已在
\qty{-1.9}{\celsius} 处达到这一凝固点：深海可降温的余地只有几度，而南极上空的空气
则没有这样的下限。
\textbf{24.} 考虑孔隙水校正后，末次盛冰期的深海比今天约 \textbf{\qty{3.3}{\celsius}
更冷}，与独立深层水估计所给出的
2 至 \qty{5}{\celsius} 相近。
\end{solution}
